\chapter{Periodic Structure of Psychological Systems}
\markboth{Periodic Structure of Psychological Systems}{Periodic Structure of Psychological Systems}

\begin{remark}[Periodicity decision for the target of this chapter (citation)]\label{rem:psych_ps_hook}
When this chapter decides whether a research target is periodic, determines its period type, or seeks its minimal positive period, reference is always made to the `Periodic Systems' chapter:
for the periodicity criterion see \cref{thm:ps_periodicity_decision},
for the discrimination of period types see \cref{def:ps_period_type}, \cref{prop:ps_type_discrimination},
for the search for the minimal positive period see \cref{def:ps_minimal_period_algorithm}, \cref{thm:ps_minimal_period_correctness},
and for the construction of period series and epicycle groups from data see \cref{thm:ps_series_reconstruction}, \cref{def:ps_epicycle_group}.
\end{remark}

\begin{remark}[Positioning of this chapter: unfolding along the five-step interface of the `Periodic Systems' chapter]\label{rem:psych_five_steps_positioning}
The periodic structure of psychological systems studies the mathematical characterisation of periodicity, rhythmicity and oscillatory phenomena in human psychological activity. This chapter does not set up a separate periodicity criterion; instead it connects psychological systems to the five-step interface specified by the `Periodic Systems' chapter:
\begin{enumerate}[label=(\arabic*)]
    \item \textbf{Decision}: the sleep--wake cycle, attention fluctuation, the emotional cycle, memory decay, subconscious processes, behavioural periods, interpersonal synchrony and other such targets are formalised as discrete or continuous dynamical systems, and \cref{thm:ps_periodicity_decision} is cited to decide whether they are periodic and to determine the period type;
    \item \textbf{Construction}: the amplitude series, the frequency series and the phase series are constructed from arousal level, emotional valence, attention resource, memory strength, subconscious activation and behavioural tendency, and assembled into the psychological epicycle group;
    \item \textbf{Adjustment}: the psychological epicycle group is dynamically adjusted under the constraints of residual energy decrease, psychological energy conservation, information conservation and attention conservation;
    \item \textbf{Convolution}: amplitude convolution, phase convolution and frequency convolution are introduced on the psychological epicycle group, and connected with autocorrelation, cross-correlation and emotional synchrony;
    \item \textbf{Integration}: the tool chains of three chapters---`Discrete Sampling and Intercalation of Periodic Signals', `Harmonic Analysis' and `Time Series Analysis'---are connected, forming the psychological periodicity diagnostic pipeline.
\end{enumerate}
The mathematical tools used in this chapter include: the two-process model, limit-cycle theory, the RC charge--discharge model, the circumplex model of emotion, the Ebbinghaus forgetting curve, the phase-oscillator model, coupled oscillation, periodic convolution, the epicycle group, the phase operator and time delay.
\end{remark}

\begin{remark}[Guide to the seven themes of this chapter]\label{rem:psych_seven_themes}
\textbf{Translation}---the approximately 24-hour main period of the sleep--wake cycle and the advance of biological rhythms through time; \textbf{Scaling}---the exponential decay of the Ebbinghaus forgetting curve and the scaling optimisation of spaced repetition; \textbf{Rotation}---the circadian rotation in the two-dimensional valence--arousal phase space and the limit cycle of bipolar disorder; \textbf{Conservation}---psychological energy conservation, information conservation and attention conservation (empirical observation / research programme, Remark~\ref{thm:psychological_conservation}); \textbf{Symmetry}---the positive--negative valence symmetry of emotion and the circadian symmetry of sleep--wake; \textbf{Closure}---the closure of the attention resource consumption--recovery cycle (Lemma~\ref{ch26_attention_cycle}); \textbf{Period}---the 90-minute sleep cycle, attention oscillation (theta/delta bands), the circadian emotional cycle and the habit loop.
\end{remark}

\section{Connection to the five-step interface of the `Periodic Systems' chapter}
\label{sec:psych_five_steps}

\subsection{Step one: periodicity decision for psychological targets}

\begin{definition}[Psychological target and decision input]\label{def:psych_target}
A \textbf{psychological target} is one of the following objects:
\begin{enumerate}[label=(\roman*)]
    \item a sleep--wake system, whose state space is the joint space of sleep pressure and circadian phase;
    \item an attention resource system, whose state space is the interval space of the attention resource $R(t)$;
    \item an emotional cycle system, whose state space is the two-dimensional (valence $V$, arousal $A$) phase space;
    \item a memory retention system, whose state space is the interval space of the memory strength $M(t)$;
    \item a behavioural period system, whose state space is the joint space of the behavioural tendency $B(t)$ and cue strength;
    \item an interpersonal emotional synchrony system, whose state space is the joint space of the emotional phases $(\theta_1,\theta_2)$ of the two parties.
\end{enumerate}
The \textbf{decision input} of a psychological target is an ordered pair
\[
\mathcal D_{\mathrm{P}}=(\mathcal S,\Phi),
\]
where $\mathcal S$ is the state space and $\Phi:\mathcal S\to\mathcal S$ is the map induced by time translation, phase rotation or feedback evolution. If $\mathcal S$ is a finite set, the target is called a \textbf{finite psychological target}; if $\mathcal S$ is an infinite set, it is called an \textbf{infinite psychological target}.
\end{definition}

\begin{definition}[Psychological decision map]\label{def:psych_periodization}
Let a psychological target be given by the decision input $\mathcal D_{\mathrm{P}}=(\mathcal S,\Phi)$. If the states of the target are generated by $\Phi$, that is, if there exists $x_0\in\mathcal S$ such that
\[
x_t=\begin{cases}\Phi^t(x_0), & t\in\mathbb Z,\\ \Phi_t(x_0), & t\in\mathbb R,\end{cases}
\]
for $t\in\mathbb Z$ the state $x_t$ is generated by iteration of $\Phi$; for $t\in\mathbb R$ it is described by the one-parameter evolution family $\{\Phi_t\}_{t\in\mathbb R}$, where each $\Phi_t:\mathcal S\to\mathcal S$, $\Phi_0=\operatorname{id}$, $\Phi_{s+t}=\Phi_s\circ\Phi_t$ for all $s,t\in\mathbb R$, and $(t,x)\mapsto\Phi_t(x)$ is continuous.
then $\Phi$ is called the \textbf{psychological decision map} of the target. The psychological target is thereby reduced to the discrete dynamical system $(\mathcal S,\Phi)$ when $t\in\mathbb Z$ and to the continuous dynamical system $(\mathcal S,\{\Phi_t\}_{t\in\mathbb R})$ when $t\in\mathbb R$; in both cases the criteria of the `Operations and Periodicity' chapter and the `Periodic Systems' chapter may be applied directly.
\end{definition}

\begin{theorem}[Periodicity decision theorem for psychological systems]\label{thm:psych_periodicity_decision}
Let the decision map of a psychological target be $\Phi$, and let $(\mathcal S,\Phi)$ be a discrete dynamical system. Then:
\begin{enumerate}[label=(\Alph*)]
    \item \textbf{Sufficient decision}: if there exists $m\in\mathbb N^+$ such that $\Phi^m=\mathrm{id}$, then the target is a uniformly periodic psychological system; every orbit is closed, and the orbit length divides $m$.
    \item \textbf{Decision equivalence}: the target is a uniformly periodic psychological system if and only if $\operatorname{ord}(\Phi)<\infty$.
    \item \textbf{Pointwise criterion}: the orbit of the target through the state $x_0$ is closed if and only if $x_0$ is a periodic point; the orbit length is then the period at that point.
    \item \textbf{Non-omissibility}: the conjunction of conserved quantities and symmetry groups cannot replace $\operatorname{ord}(\Phi)<\infty$; a quasi-periodic psychological target does not belong to the periodic psychological systems.
\end{enumerate}
\end{theorem}

\begin{proof}
The \cref{thm:psych_periodicity_decision} of this chapter is the domain instantiation of \cref{thm:ps_periodicity_decision} in the `Periodic Systems' chapter:
it suffices to replace the decision map $f$ and the evolution group by the $\Phi$ of this chapter, and the four items---(A) sufficient decision, (B) decision equivalence, (C) pointwise criterion and
(D) non-omissibility---hold one by one. By the convention of \cref{rem:ps_applications_hook}, the argument is not repeated here.
\end{proof}

\begin{corollary}[Types of psychological period]\label{cor:psych_period_type}
After the decision by \cref{thm:psych_periodicity_decision}, a psychological target falls into one of the following types:
\begin{enumerate}[label=(\roman*)]
    \item \textbf{Trivially periodic type}: the least period is $m_0=1$, and the state is a fixed point;
    \item \textbf{Strictly periodic type}: $1<m_0<\infty$, and the state is strictly periodic;
    \item \textbf{Uniformly periodic type}: $\operatorname{ord}(\Phi)=m<\infty$, and all orbit lengths divide $m$;
    \item \textbf{Quasi-periodic type}: $\operatorname{ord}(\Phi)=\infty$; the orbits are not closed but are dense on a compact carrier structure;
    \item \textbf{Eventually periodic type}: there exists $N$ such that the target is periodic for $t\ge N$;
    \item \textbf{Discretely periodic type}: the carrier structure is $\mathbb Z$, and the normalised period ratio is $\rho=T_x/\Delta x\in\mathbb Q$.
    \item \textbf{Asymptotically periodic (damped) type}: the orbit is not closed, but tends asymptotically to a closed orbit or to a fixed point, the residual decaying as $e^{-\kappa t}$ (see \cref{def:ps_period_type}(vii)).
\end{enumerate}
Here (i)--(iii) and (vi) belong to the periodic psychological systems; (iv) does not belong to the periodic psychological systems but can be handled with bounded-approximation tools; (v) is a periodic psychological system only in its stable stage; (vii), like (iv), is a type in the limiting sense; it is not a periodic psychological system and can be handled with bounded-approximation tools only in its stable stage.
\end{corollary}

\begin{proof}
Each item follows by substituting the psychological target into \cref{thm:psych_periodicity_decision} and \cref{def:ps_period_type} of the `Periodic Systems' chapter.
\end{proof}

\begin{lemma}[Decision for the sleep--wake cycle]\label{lem:sleep_periodic}
Let the sleep--wake system be described by the two-process model: process S accumulates linearly during the waking period and is cleared exponentially during the sleep period; process C oscillates with a period of about 24 hours. Then the sleep--wake system is a uniformly periodic system, whose main period is
\[
T_{\text{sleep-wake}}\approx 24\ \text{hours},
\]
and an ultradian NREM--REM cycle of about $90$ minutes is nested within this main period.
\end{lemma}

\begin{proof}
By Proposition~\ref{prop:two_process_sleep}, the sleep propensity $P(t)=S(t)-C(t)$ crosses the threshold in the evening and at night, triggering sleep, and falls back below the threshold in the early morning, triggering waking. Process C oscillates with a period of 24 hours, and process S accumulates and is cleared periodically in the wake--sleep cycle; the superposition of the two gives the sleep--wake system a main period of about 24 hours. Hence the decision map $\Phi$ satisfies $\Phi^{T_{\text{sleep-wake}}}=\mathrm{id}$, and the system is uniformly periodic.
\end{proof}

\begin{lemma}[Periodicity decision for attention fluctuation]\label{lem:attention_periodic}
Let the attention resource $R(t)$ be consumed at rate $\alpha$ during task periods and restored at rate $\beta$ during rest periods. Then the consumption--recovery cycle of the attention resource is a uniformly periodic system, whose period is given by Lemma~\ref{ch26_attention_cycle}.
\end{lemma}

\begin{proof}
In the steady-state cycle of alternating task and rest, the attention resource charges and discharges periodically within $[R_{\min},R_{\max}]$. By the closed-form solution of Lemma~\ref{ch26_attention_cycle}, $R_{\min}$ and $R_{\max}$ are uniquely determined by the closure conditions, and the cycle repeats with period $T=t_1+t_2$. Hence the decision map $\Phi$ satisfies $\Phi^T=\mathrm{id}$, and the system is uniformly periodic.
\end{proof}

\begin{lemma}[Decision for the emotional cycle]\label{lem:emotion_periodic}
Let the emotional state rotate about the equilibrium point in the two-dimensional (valence $V$, arousal $A$) phase space with angular frequency $\omega$, with return rate $\mu$. Then:
\begin{enumerate}[label=(\roman*)]
    \item if $\mu>0$ and the oscillation persists (external driving), the emotional system is a uniformly periodic system;
    \item if $\mu>0$ and there is no external driving, the emotional orbit is a decaying spiral and belongs to the asymptotically periodic (damped) type;
    \item if $\mu\approx 0$ (the limit cycle of bipolar disorder), the emotional system is a uniformly periodic system, with amplitude determined by the pathological limit cycle.
\end{enumerate}
\end{lemma}

\begin{proof}
By the proof of Proposition~\ref{prop:bipolar_cycle}, the linearisation of the emotional dynamics near the equilibrium point is $\dot X=JX$, with eigenvalues $\lambda_\pm=-\mu\pm i\omega$. When $\mu>0$ the orbit is a decaying spiral and belongs to the asymptotically periodic (damped) type; when $\mu\approx 0$ the orbit is approximately closed (uniformly periodic type); when there is persistent external driving (such as circadian driving), the steady-state response is of uniformly periodic type.
\end{proof}

\begin{lemma}[Periodicity decision for the forgetting curve]\label{lem:memory_periodic}
Let the memory strength $M(t)$ decay exponentially as $M(t)=M_0 e^{-t/\tau}$ in the absence of review, and be reset to $M_0$ on review. Then the review--forgetting cycle is a uniformly periodic system, whose period is the sum of the review intervals.
\end{lemma}

\begin{proof}
Each review resets $M$ to $M_0$; between two reviews, $M$ decays exponentially to the threshold $M_{\min}$. By Proposition~\ref{prop:spaced_repetition}, the optimal time interval before the $n$-th review is $\Delta t_n$. When the sequence of review intervals is periodic, the review--forgetting cycle constitutes a uniformly periodic system.
\end{proof}

\begin{lemma}[Periodicity decision for habit formation]\label{lem:habit_periodic}
Let the behavioural tendency $B(t)$ converge to the asymptotic value $B_\infty$ at rate $\gamma$ under repeated cues. Then the process of habit formation is of asymptotically periodic (damped) type---it converges to a steady state after the time scale $\tau_{\text{habit}}=1/\gamma$. If the cues themselves are periodic, $B$ exhibits a periodic response.
\end{lemma}

\begin{proof}
By Proposition~\ref{prop:habit_formation}, under repeated cues $B(t)=B_\infty(1-e^{-\gamma t})$, converging to $B_\infty$. If the cues themselves repeat with period $T_{\text{cue}}$, then in the steady state $B$ exhibits a periodic response of period $T_{\text{cue}}$. The two cases correspond to the asymptotically periodic (damped) type and the uniformly periodic type respectively.
\end{proof}

\begin{definition}[Least positive period algorithm for psychological systems]\label{def:psych_minimal_period_algorithm}
This definition has the same content as \cref{def:ps_minimal_period_algorithm} in the section `Periodicity Decision: Is the Target Periodic?' (\cref{sec:ps_periodicity_decision}) of Chapter~\ref{chap:periodic_system}; to avoid repeating the statement, only the reference anchor is retained here, and the complete statement is given there. The only difference here is terminology (psychological target / psychological sample); the algorithmic steps and the conclusions are exactly the same.
\end{definition}

\begin{theorem}[Correctness of the least positive period algorithm for psychological systems]\label{thm:psych_minimal_period_correctness}
This theorem has the same content as \cref{thm:ps_minimal_period_correctness} in the section `Periodicity Decision: Is the Target Periodic?' (\cref{sec:ps_periodicity_decision}) of Chapter~\ref{chap:periodic_system}; to avoid repeating the statement, only the reference anchor is retained here, and the complete statement and proof are given there. The only difference here is terminology (psychological target / psychological sample); the algorithmic steps and the conclusions are exactly the same.
\end{theorem}

\begin{example}[Examples of periodicity decisions for psychological systems]\label{ex:psych_periodic_examples}
\begin{enumerate}[label=(\roman*)]
    \item sleep--wake: $T_{\text{sleep-wake}}\approx 24$ h, uniformly periodic type;
    \item ultradian rhythm within sleep: $T_{\text{ultradian}}\approx 90$ min, uniformly periodic type;
    \item attention fluctuation: $T_{\text{attention}}=t_1+t_2$, uniformly periodic type;
    \item circadian emotion: $T_{\text{emotion}}=24$ h (with circadian driving), uniformly periodic type;
    \item bipolar disorder: $T_{\text{bipolar}}$ of weeks to months, uniformly periodic type (limit cycle);
    \item memory review: $T_{\text{review}}=\sum_n \Delta t_n$, uniformly periodic type (under periodic review);
    \item habit formation: asymptotically periodic (damped) type (without periodic cues) or uniformly periodic type (with periodic cues).
\end{enumerate}
\end{example}

\begin{remark}[Conserved quantities do not guarantee periodicity]\label{rem:psych_invariant_not_period}
Energy conservation, information conservation and attention conservation in psychological systems are auxiliary conditions for periodic structure, but they do not constitute sufficient conditions. For example, a mental state that possesses energy conservation but lacks a feedback loop will not spontaneously produce periodic oscillation. Psychological periodic structure must be decided by \cref{thm:psych_periodicity_decision} or by the lemmas above.
\end{remark}

\subsection{Step two: construction of the psychological period series and the epicycle group}

\begin{definition}[Psychological amplitude, frequency and phase series]\label{def:psych_series}
Let the state or output sequence of a psychological target be $\{z_t\}$. If it can be written as
\[
z_t=\sum_k A_k e^{i(\omega_k t+\phi_k)}+\varepsilon_t,
\]
then
\[
\mathcal A_z(t)=\sum_k A_k e^{ikt},\qquad
\mathcal F_z(s)=\sum_k \omega_k s^k,\qquad
\mathcal P_z(t)=\sum_k e^{i(kt+\phi_k)}
\]
are called respectively the \textbf{amplitude series}, the \textbf{frequency series} and the \textbf{phase series} of the psychological target. The three agree with \cref{def:amplitude_series}, \cref{def:frequency_series} and \cref{def:phase_series} of the `Series and Convolution' chapter.
\end{definition}

\begin{definition}[Psychological epicycle group]\label{def:psych_epicycle_group}
Let the period series of a psychological target be given by the amplitude series $\{A_k\}$, the frequency series $\{\omega_k\}$ and the phase series $\{\phi_k\}$. For each $k$ define the epicycle
\[
E_k:\theta\mapsto A_k e^{i(\theta+\phi_k)},\qquad \theta=\omega_k t,
\]
whose collection
\[
\mathfrak E_{\mathrm{P}}=\{(A_k,\omega_k,\phi_k)\}_{k\in\mathbb Z}
\]
is called the \textbf{epicycle group} of the psychological target. Its superposition is the psychological period series
\[
z_t=\sum_k E_k(\omega_k t)=\sum_k A_k e^{i(\omega_k t+\phi_k)}.
\]
\end{definition}

\begin{theorem}[Construction of the period series from psychological data]\label{thm:psych_series_reconstruction}
Let the output of a psychological target be a finite periodic sequence of period $N$. Then:
\begin{enumerate}[label=(\roman*)]
    \item the output can be expanded as a discrete Fourier series
    \[
    z_t=\sum_{k=0}^{N-1}c_k e^{i2\pi kt/N},\qquad
    c_k=\frac1N\sum_{t=0}^{N-1}z_t e^{-i2\pi kt/N};
    \]
    \item the amplitude series is given by $|c_k|$, the phase series by $\arg c_k$ and the frequency series by $2\pi k/N$;
    \item if the output is a sleep--wake sequence, the fundamental frequency corresponds to the main period of about 24 hours, and the harmonics correspond to the ultradian rhythm of about 90 minutes and the semi-diurnal rhythm of about 12 hours; if the output is an attention resource sequence, the fundamental frequency corresponds to the period of the task--rest cycle.
\end{enumerate}
\end{theorem}

\begin{proof}
(i) The discrete Fourier series of a finite periodic sequence is a standard result: the basis functions $\{e^{i2\pi kt/N}\}_{k=0}^{N-1}$ form an orthogonal basis of $\mathbb C^N$.
(ii) This follows directly from $c_k=A_k e^{i\phi_k}$.
(iii) For the sleep--wake sequence with $N=24$ (hourly sampling), the fundamental frequency corresponds to the 24-hour main period, and the $k=16$ harmonic corresponds to the 90-minute ultradian period ($24\times 60/90=16$); for the attention sequence, the fundamental frequency corresponds to the period of the task--rest cycle.
\end{proof}

\begin{proposition}[Three-parameter correspondence: radius, phase and speed]\label{prop:psych_epicycle_correspondence}
The three parameters of the psychological epicycle group $\mathfrak E_{\mathrm{P}}=\{(A_k,\omega_k,\phi_k)\}_{k}$ correspond one-to-one to three geometric--psychological quantities:
\begin{enumerate}[label=(\roman*)]
    \item \textbf{Radius corresponds to amplitude}: the radius $|E_k|=A_k$ of $E_k$ is the amplitude of that psychological oscillation component;
    \item \textbf{Phase corresponds to rotational delay}: the initial phase $\phi_k$ corresponds to the time translation $\tau_k=\phi_k/\omega_k$; when $\phi_k>0$ this component \textbf{leads} the baseline by $\tau_k$ (equivalently, the component coincides with the baseline once the baseline is postponed by $\tau_k$);
    \item \textbf{Frequency corresponds to rotational speed}: $\omega_k$ is the angular frequency of that psychological oscillation component.
\end{enumerate}
\end{proposition}

\begin{proof}
The \cref{prop:psych_epicycle_correspondence} of this chapter corresponds exactly to \cref{prop:ps_epicycle_correspondence} in the `Periodic Systems' chapter,
and the three correspondences---radius to amplitude, phase to rotational delay, and frequency to rotational speed---are obtained by translating each of them to the epicycle group of this chapter.
The proof is given there: (i) from $|E_k|=A_k$ and the energy decomposition, (ii) from the phase-translation theorem, (iii) from the definition of the frequency series.
\end{proof}

\begin{example}[The period series of sleep--wake]\label{ex:psych_sleep_series}
The sleep--wake sequence contains a main period of about 24 hours, an ultradian period of about 90 minutes and a semi-diurnal period of about 12 hours. In the psychological epicycle group, the three appear as epicycle components of different radii and different angular frequencies---the main-period amplitude is the largest, whilst the ultradian and semi-diurnal components have smaller amplitudes.
\end{example}

\begin{remark}[The period series is a complex function]\label{rem:psych_series_complex}
Let $z=e^{i\omega_0t}$. Then the psychological period series
\[
z_t=\sum_k A_k e^{i(\omega_k t+\phi_k)}
\]
when all components are commensurable with $\omega_0$ (that is, $\omega_k=k\omega_0$, $k\in\mathbb Z$) it can be written as the Laurent series $\sum_k c_k z^k$, where $c_k=A_k e^{i\phi_k}$: in that case the psychological period series is a complex function on the unit circle, and otherwise it is merely a function on the torus.
\end{remark}

\subsection{Step three: dynamic adjustment of the psychological epicycle group}

\begin{definition}[Dynamic adjustment of the psychological epicycle group]\label{def:psych_adaptive_adjustment}
Let the current psychological epicycle group be $\mathfrak E^{(j)}$, with partial sum
\[
z^{(j)}(t)=\sum_{k\in\mathfrak E^{(j)}}A_k e^{i(\omega_k t+\phi_k)},
\]
and residual $r^{(j)}=z-z^{(j)}$. \textbf{Dynamic adjustment} consists of two rules:
\begin{enumerate}[label=(\roman*)]
    \item \textbf{Cardinality adjustment}: if the residual energy $\|r^{(j)}\|^2>\delta$, a new epicycle is added according to the peak of the residual periodogram;
    \item \textbf{Parameter adjustment}: the parameters of the existing epicycles are updated by regularised least squares or under the conservation constraints.
\end{enumerate}
In psychological systems, cardinality adjustment corresponds to adding a biological oscillatory mode (such as a newly discovered rhythm band); parameter adjustment corresponds to tuning the amplitude, phase and frequency so that the psychological energy, information and attention constraints are satisfied.
\end{definition}

\begin{proposition}[Monotonicity of cardinality adjustment]\label{prop:psych_adaptive_growth}
After a dominant-frequency component of the residual has been added as in \cref{def:psych_adaptive_adjustment}(i), the residual energy is monotonically non-increasing:
\[
\|r^{(j+1)}\|^2\le\|r^{(j)}\|^2,
\]
with equality if and only if the residual has no energy at the selected frequency.
\end{proposition}

\begin{proof}
The \cref{prop:psych_adaptive_growth} of this chapter is consistent with the dynamic adjustment mechanism of \cref{def:ps_adaptive_adjustment} in the `Periodic Systems' chapter:
adding an epicycle removes the component of $r^{(j)}$ along $\hat\omega$ from the residual, and the inequality follows from the orthogonal decomposition
$\|r^{(j)}\|^2=\|r^{(j+1)}\|^2+\|\text{the removed component}\|^2$,
with equality if and only if the norm of the removed component is zero. The complete argument is given there.
\end{proof}

\begin{theorem}[Convergence of dynamic adjustment]\label{thm:psych_adaptive_convergence}
Let the psychological target be a finite periodic sequence of period $T_0$ with $\delta=0$. Then after at most $\lfloor T_0/2\rfloor+1$ steps of dynamic adjustment the residual is zero, and the resulting psychological epicycle group is the exact Fourier data of the target.
\end{theorem}

\begin{proof}
The discrete Fourier series of a sequence of period $T_0$ is complete on $\lfloor T_0/2\rfloor+1$ harmonics. At each step a dominant-frequency component of the residual is extracted as in \cref{prop:psych_adaptive_growth}, and the residual energy decreases strictly; hence after finitely many steps the residual is zero.
\end{proof}

\begin{definition}[Structural metric adjustment]\label{def:psych_structure_metric_adjustment}
Let the structural metric of a psychological target be
\[
\mathcal I=(\mathcal I_{\text{energy}},\mathcal I_{\text{information}},\mathcal I_{\text{attention}})
=(\text{psychological energy conservation degree},\ \text{information conservation degree},\ \text{attention conservation degree}).
\]
\textbf{Structural-metric adjustment} consists in choosing $\theta$ in the parameter space $\Theta$: at each step take a neighbouring parameter for which $\mathcal I(\theta)$ is \textbf{strictly decreasing} in the lexicographic order (minimisation); if no such parameter exists, the process stops, and the parameter is then locally optimal.
\end{definition}

\begin{proposition}[Finite convergence of structural metric adjustment]\label{prop:psych_structure_metric_convergence}
Let the parameter space $\Theta$ be finite, and let each adjustment make $\mathcal I$ strictly decreasing in the lexicographic order. Then the adjustment terminates after finitely many steps at a locally optimal parameter.
\end{proposition}

\begin{proof}
Since $\Theta$ is finite, $\mathcal I$ takes only finitely many values, and the lexicographic order is a total order on a finite set; a strictly decreasing sequence in that order cannot continue indefinitely, so the adjustment terminates after finitely many steps. At termination there is no neighbouring parameter that makes $\mathcal I$ strictly decrease in the lexicographic order, that is, the parameter is locally optimal.
\end{proof}

\begin{remark}[Boundaries of dynamic adjustment]\label{rem:psych_adaptive_boundary}
\cref{thm:psych_adaptive_convergence} relies on finite periodic sequences. For continuous spectra or quasi-periodic psychological targets there is no exact representation by a finite epicycle group, and only bounded approximation is possible. Structural metric adjustment does not guarantee global optimality.
\end{remark}

\subsection{Step four: psychological convolution}

\begin{definition}[Psychological amplitude, phase and frequency convolutions]\label{def:psych_convolution_types}
Let two psychological epicycle groups be
\[
\mathfrak E=\{(A_n,\omega_n,\phi_n)\},\qquad
\mathfrak E'=\{(B_m,\nu_m,\psi_m)\}.
\]
Define:
\begin{enumerate}[label=(\roman*)]
    \item \textbf{Amplitude convolution}: $\mathbb A_{mn}=A_nB_m$;
    \item \textbf{Phase convolution}: $\mathbb P_{mn}=e^{i(\phi_n-\psi_m)}$;
    \item \textbf{Frequency convolution}: $\mathbb F_{mn}=\omega_n\nu_m$;
    \item \textbf{Trend convolution}: $\mathbb T_{mn}=A_nB_m e^{i(\phi_n-\psi_m)}=A_nB_m\mathbb P_{mn}$.
\end{enumerate}
These agree with \cref{def:amplitude_convolution}, \cref{def:phase_convolution}, \cref{def:frequency_convolution} and \cref{def:tendency_convolution} of the `Series and Convolution' chapter.
\end{definition}

\begin{theorem}[Action of the three convolutions on psychological epicycle groups]\label{thm:psych_convolution_action}
Let the epicycle groups of two psychological targets be $\mathfrak E$ and $\mathfrak E'$ respectively. Then:
\begin{enumerate}[label=(\roman*)]
    \item amplitude convolution acts only on the radius;
    \item phase convolution acts only on the phase;
    \item frequency convolution acts only on the rotational speed;
    \item trend convolution preserves both amplitude and phase.
\end{enumerate}
\end{theorem}

\begin{proof}
This follows directly from \cref{def:psych_convolution_types}.
\end{proof}

\begin{proposition}[Autocorrelation and trend convolution]\label{prop:psych_autocorr_phase_convolution}
Let $\{z_t\}$ be a psychological period sequence, with autocorrelation function
\[
R_z(h)=\sum_t z_t\overline{z_{t+h}}.
\]
In the language of the psychological epicycle group, the autocorrelation function is the projection of the trend convolution onto the diagonal:
\[
R_z(h)=\sum_n\mathbb T_{nn}e^{-i\omega_n h}.
\]
\end{proposition}

\begin{proof}
The \cref{prop:psych_autocorr_phase_convolution} of this chapter is the realisation in this domain of the convolution-action conclusion of \cref{thm:ps_convolution_action} in the `Periodic Systems' chapter:
substituting $z_t=\sum_n A_n e^{i(\omega_n t+\phi_n)}$ into the definition of the autocorrelation, cancelling all off-diagonal cross terms by orthogonality and retaining only the diagonal term of the trend convolution
(the factorisation of the three convolutions, \cref{prop:ps_convolution_factorization}), which yields $R_z(h)=\sum_n\mathbb T_{nn}e^{-i\omega_n h}$.
The complete expansion is given there.
\end{proof}

\begin{remark}[Relation between psychological convolution and periodic convolution]\label{rem:psych_convolution_relation}
Psychological convolution is not a new convolution operation; it is an instance on the psychological epicycle group of the periodic convolution of the `Series and Convolution' chapter. Autocorrelation is the projection of the trend convolution onto the diagonal; cross-correlation and emotional synchrony are respectively the concrete manifestations in psychological science of amplitude convolution and frequency convolution.
\end{remark}

\subsection{Step five: integration with discrete sampling, harmonic analysis and time series analysis}

\begin{definition}[The psychological periodicity diagnostic pipeline]\label{def:psych_periodic_diagnosis}
For a given psychological system, diagnosis is carried out using the framework of periodic mathematics:
\begin{enumerate}[label=(\roman*)]
    \item \textbf{Periodicity decision}: verify the period type of the psychological target (\cref{thm:psych_periodicity_decision});
    \item \textbf{Series construction}: construct the amplitude, frequency and phase series and the epicycle group from arousal, emotion, attention and memory data (\cref{thm:psych_series_reconstruction});
    \item \textbf{Conservation diagnosis}: check the conservation of psychological energy, information and attention as an empirical observation and a research programme (see Remark \cref{thm:psychological_conservation});
    \item \textbf{Convolution diagnosis}: compute the autocorrelation (the diagonal projection of the trend convolution, \cref{prop:psych_autocorr_phase_convolution}) and the phase convolution;
    \item \textbf{Integration diagnosis}: connect to the tool chains of discrete sampling, harmonic analysis and time series analysis.
\end{enumerate}
\end{definition}

\begin{theorem}[The psychological periodicity diagnostic pipeline and its interface with three chapters]\label{thm:psych_diagnosis_integration}
The five steps of \cref{def:psych_periodic_diagnosis} connect respectively to:
\begin{enumerate}[label=(\roman*)]
    \item the `Discrete Sampling and Intercalation of Periodic Signals' chapter: \cref{def:uniform_sampling_parallel}, \cref{thm:nyquist_shannon_parallel}, \cref{thm:intercalation_leakage_suppression};
    \item the `Harmonic Analysis' chapter: \cref{def:fourier_series}, \cref{thm:parseval_fourier}, \cref{prop:twelve_to_six};
    \item the `Time Series Analysis' chapter: \cref{def:ts_period_detection}, \cref{prop:ts_period_significance}, \cref{def:ts_window_validation}.
\end{enumerate}
Hence the psychological periodicity diagnosis is not an independent method, but an application of the unified tool chain of periodic mathematics to psychological systems.
\end{theorem}

\begin{proof}
The \cref{thm:psych_diagnosis_integration} of this chapter is the implementation in this domain of the five-step interface of the `Periodic Systems' chapter (\cref{tab:ps_five_steps}):
the periodicity decision uses sampling and spectral estimation, the frequency decomposition and phase kernels use the `Harmonic Analysis' chapter, and the periodogram peaks, the Fisher test and the
window-function validation use the `Time Series Analysis' chapter. Hence the diagnostic pipeline is synthesised from the tool chains of the three chapters; the proof is given there.
\end{proof}

\begin{remark}[The order of operations after integration]\label{rem:psych_integration_order}
The recommended order for practical diagnosis is: first decide whether the psychological target is periodic (\cref{thm:psych_periodicity_decision}); then construct the Fourier data of the psychological data and the epicycle group (\cref{thm:psych_series_reconstruction}); then optimise the parameters using the conservation constraints and dynamic adjustment (\cref{thm:psych_adaptive_convergence}, \cref{prop:psych_structure_metric_convergence}); then carry out the psychological convolution and autocorrelation analysis (\cref{prop:psych_autocorr_phase_convolution}); and finally output the diagnosis according to \cref{def:psych_periodic_diagnosis}.
\end{remark}

\section{Mathematical foundations of psychological systems}

\begin{definition}[Basic quantities of psychological periodic systems]\label{def:psychological_quantities}
A psychological periodic system is described by the following basic variables:
\begin{enumerate}[label=(\roman*)]
    \item \textbf{Arousal level} $A(t) \in [0,1]$: the degree of alertness of the individual at time $t$. $A \to 1$ indicates high arousal (tension, excitement), and $A \to 0$ low arousal (drowsiness, relaxation). The arousal level fluctuates on a period of about 24 hours and also exhibits a rhythm on the ultradian period of about 90 minutes.
    \item \textbf{Emotional valence} $V(t) \in [-1,1]$: the degree of subjective pleasantness of the individual at time $t$. $V > 0$ indicates positive emotion (pleasantness), $V < 0$ negative emotion (unpleasantness) and $V = 0$ neutrality. Emotional valence and arousal level constitute the two-dimensional circumplex model of emotion.
    \item \textbf{Attention resource} $R(t) \in [0,1]$: the total amount of cognitive resources available to the individual at time $t$. $R \to 1$ indicates abundant attention and $R \to 0$ attentional depletion. The attention resource is consumed during task performance and restored during rest---the alternation of consumption and recovery constitutes the cycle of attention.
    \item \textbf{Memory strength} $M(t) \in [0,1]$: the retrieval probability of a particular memory trace at time $t$. $M \to 1$ indicates vivid memory and $M \to 0$ forgetting. In the absence of review, memory strength decays approximately exponentially.
    \item \textbf{Subconscious activation} $U(t) \in [0,1]$: the degree of activity of the psychological processes running below the threshold of consciousness. $U \to 1$ indicates that subconscious processes are highly active (such as memory consolidation during sleep, or the automatic processing involved in habituation), and $U \to 0$ that subconscious processes are in a resting state.
    \item \textbf{Behavioural tendency} $B(t) \in [0,1]$: the probability that the individual exhibits a particular behavioural pattern at time $t$. The behavioural tendency is determined jointly by psychological state (conscious and subconscious) and environmental stimuli.
\end{enumerate}
The core feature of psychological systems is \textbf{multiscale nested periodicity}---rhythms on different time scales (neural oscillations at the millisecond scale, attention fluctuations at the second scale, emotional changes at the hour scale, the circadian sleep cycle, the interannual behavioural period) coexist in one and the same psychological system, and span both the conscious and the subconscious level.
\end{definition}

\section{The sleep--wake cycle}

\begin{definition}[Periodic structure of sleep]\label{def:sleep_cycle}
Human sleep consists of several \textbf{sleep cycles} of about 90 minutes. Each sleep cycle contains:
\begin{enumerate}[label=(\roman*)]
    \item \textbf{NREM sleep} (non-rapid eye movement sleep): divided into N1 (sleep onset), N2 (light sleep) and N3 (deep sleep / slow-wave sleep). Brain-wave frequency gradually decreases and amplitude gradually increases---$\delta$ waves ($0.5$--$4$ Hz) dominate in the N3 stage.
    \item \textbf{REM sleep} (rapid eye movement sleep): brain waves show a high-frequency, low-amplitude desynchronised pattern (similar to the waking state), the eyes move rapidly, the body's muscles are paralysed, and dreams occur mainly in this stage.
\end{enumerate}
A whole night's sleep usually contains $4$--$6$ sleep cycles. The first half of the night is dominated by deep sleep (N3) and the second half by REM sleep---the duration of REM sleep grows from about $10$ minutes in the first cycle to about $60$ minutes in the last cycle.

The approximately 24-hour main sleep--wake period is the internalisation of the Earth's rotation period. The approximately 90-minute sleep cycle is an ultradian rhythm nested within the main period---the alternation of NREM and REM appears as a nested period in the structure of sleep. Deep sleep (NREM) corresponds to slow-wave oscillation (low frequency, high amplitude, an outer large epicycle), and REM sleep to fast-wave oscillation (high frequency, low amplitude, inner small epicycles). The increase in the proportion of REM sleep over a whole night's sleep is a phase drift within the nested period---the REM component gradually strengthens through the night.
\end{definition}

\begin{proposition}[Two-process model of sleep]\label{prop:two_process_sleep}
The regulation of sleep consists of two independent processes:
\begin{enumerate}[label=(\roman*)]
    \item \textbf{Process S (sleep pressure)}: during waking, sleep-promoting substances such as adenosine accumulate in the brain, and sleep pressure increases linearly. During sleep, these substances are cleared, and sleep pressure decays exponentially. Process S manifests as the linear accumulation of sleep pressure during the waking period and its exponential clearance during the sleep period---analogous to the charging (waking period) and discharging (sleep period) of a capacitor.
    \item \textbf{Process C (circadian rhythm)}: an approximately 24-hour circadian signal generated by the suprachiasmatic nucleus (SCN) that modulates sleep propensity. Process C promotes sleep at night and waking during the day. Process C is a periodic drive---SCN neurons oscillate with a period of about 24 hours.
\end{enumerate}
The superposition of the two processes determines the timing and the depth of sleep: sleep occurs at the moment when process S is at its highest and process C most promotes sleep (usually at night). The two-process model of sleep is an instance in which a linear accumulation (process S) and a periodic oscillation (process C) couple to produce the sleep--wake cycle.
\end{proposition}

\begin{proof}
Write the waking period as $t \in [0, T_w)$ and the sleep period as $[T_w, T_w + T_s)$ ($T_w + T_s = 24$ h). The linearised model of process S is: during the waking period $S(t) = S_{\min} + \rho_w t$ ($\rho_w > 0$ is the rate of accumulation of sleep pressure), and during the sleep period $S(t) = S_{\max} e^{-t/\tau_s}$ ($\tau_s$ is the clearance time constant). Process C is approximated by a sinusoid $C(t) = C_0 \cos\!\big(\tfrac{2\pi t}{24}\big)$ (negative at night, which promotes sleep). The sleep propensity is $P(t) = S(t) - C(t)$: waking $\Leftrightarrow$ $P(t) < \theta$, sleep $\Leftrightarrow$ $P(t) \ge \theta$ (threshold $\theta$). Since $S$ increases monotonically during the waking period and $C$ decreases monotonically in the evening, $P$ necessarily crosses the threshold in the evening and at night---sleep is triggered in a `gated' manner during the interval in which $S$ is high and $C$ is low; in the early morning $C$ rises again and $S$ has already been cleared, so $P$ falls back below the threshold---waking. The superposition of the two processes thus rigorously yields the approximately 24-hour sleep--wake cycle.
\end{proof}

\section{Periodic fluctuation of attention}

\begin{definition}[Oscillation of attention]\label{def:attention_oscillation}
Human attention is not continuous and uniform; rather, it \textbf{fluctuates rhythmically} at several frequencies:
\begin{enumerate}[label=(\roman*)]
    \item \textbf{High-frequency attention oscillation} ($4$--$8$ Hz, theta band): attention fluctuates on a time scale of about $125$--$250$ milliseconds. This corresponds to periodic variation in the perceptual sampling rate---the `bit width' of attention (the number of targets that can be attended to simultaneously) scales at the theta frequency.
    \item \textbf{Mid-frequency attention oscillation} ($0.5$--$2$ Hz, delta band): attention fluctuates on a time scale of about $0.5$--$2$ seconds. This corresponds to switching of the focus of attention among several targets---the `bandwidth' of attention (the temporal resolution of attentional switching) varies at the delta frequency.
    \item \textbf{Low-frequency attention oscillation} ($<0.1$ Hz): attention fluctuates on a time scale of minutes. This corresponds to fluctuations in overall cognitive resources---task performance exhibits a periodic fluctuation of about $10$--$20$ minutes.
\end{enumerate}
The periodic fluctuation of attention manifests as an alternation of focusing and dispersion: when focusing, cognitive resources are concentrated on a single target; when dispersing, resources are distributed among several targets. The focusing and dispersion of attention determine the effective flux of information from the external world to the conscious mind.
\end{definition}

\begin{proposition}[Model of attention resource consumption and recovery]\label{prop:attention_resource}
Let the attention resource $R(t)$ be consumed at rate $\alpha$ during task performance and restored at rate $\beta$ during rest:
\[
\frac{dR}{dt} = -\alpha \cdot \mathbf{1}_{\text{task}}(t) \cdot R(t) + \beta \cdot \mathbf{1}_{\text{rest}}(t) \cdot (1 - R(t)),
\]
where $\mathbf{1}_{\text{task}}(t)$ and $\mathbf{1}_{\text{rest}}(t)$ are indicator functions. The alternation of consumption and recovery produces the periodic fluctuation of the attention resource.

The consumption--recovery cycle of the attention resource produces periodic behaviour through linear consumption and exponential recovery: as in the charge--discharge process of an RC circuit, the attention resource decays exponentially during task periods and is restored exponentially during rest periods. The time constants $\tau_{\text{consumption}} = 1/\alpha$ and $\tau_{\text{recovery}} = 1/\beta$ determine the duration of the attention cycle. In prolonged tasks, the optimisation of rest intervals is equivalent to finding the best attention cycle---the one that maximises the average attention resource during the task.
\end{proposition}

\begin{lemma}[Closed-form solution for the attention charge--discharge cycle]\label{ch26_attention_cycle}
Let the task period have duration $t_1$ (consumption rate $\alpha$) and the rest period duration $t_2$ (recovery rate $\beta$), the cycle repeating with period $T = t_1 + t_2$. Then the steady-state cycle satisfies:
\[
R_{\min} = \frac{(1 - e^{-\beta t_2})\, e^{-\alpha t_1}}{1 - e^{-\alpha t_1} e^{-\beta t_2}}, \qquad
R_{\max} = \frac{1 - e^{-\beta t_2}}{1 - e^{-\alpha t_1} e^{-\beta t_2}},
\]
and the period-average resource is
\[
\bar R = \frac{1}{T}\Big( \frac{R_{\max} - R_{\min}}{\alpha} + t_2 - \frac{(1 - R_{\max})(e^{\beta t_2} - 1)}{\beta} \Big).
\]
\end{lemma}

\begin{proof}
The solution during the task period is $R(t) = R_{\max} e^{-\alpha t}$, and during the rest period $R(t) = 1 - (1 - R_{\min}) e^{-\beta t}$. The period closure conditions are:
\[
\begin{aligned}
R_{\min} &= R_{\max} e^{-\alpha t_1}, \\
R_{\max} &= 1 - (1 - R_{\min}) e^{-\beta t_2}.
\end{aligned}
\]
Solving these simultaneously gives $R_{\min}$ and $R_{\max}$ as above. The average resource is computed as $\bar R = \frac{1}{T}\big(\int_0^{t_1} R_{\max}e^{-\alpha t} dt + \int_0^{t_2} [1-(1-R_{\min})e^{-\beta t}] dt\big)$: the first term is $= (R_{\max}-R_{\min})/\alpha$; the second term is $= t_2 - (1-R_{\min})(1-e^{-\beta t_2})/\beta$, which by the closure condition $1 - R_{\min} = (1-R_{\max}) e^{\beta t_2}$ becomes $t_2 - (1-R_{\max})(e^{\beta t_2}-1)/\beta$.
\end{proof}

\begin{example}[The attention cycle of the Pomodoro Technique]\label{ch26_pomodoro_example}
Take $\alpha = 0.10$/min (high-intensity task), $\beta = 0.25$/min (recovery during a short rest), $t_1 = 25$ min and $t_2 = 5$ min. Substituting into Lemma~\ref{ch26_attention_cycle}: $e^{-\alpha t_1} = e^{-2.5} \approx 0.082$, $e^{-\beta t_2} = e^{-1.25} \approx 0.287$. This gives
\[
R_{\min} = \frac{(1-0.287)\times 0.082}{1 - 0.082\times 0.287} \approx \frac{0.0584}{0.9765} \approx 0.060, \qquad
R_{\max} = \frac{0.713}{0.9765} \approx 0.730.
\]
That is, within each Pomodoro cycle the attention resource charges and discharges between $6\%$ and $73\%$, with period average $\bar R \approx 0.30$. If the rest is removed ($t_2 = 0$), then $R_{\min} \to 0$ and $\bar R$ drops markedly---the presence of rest intervals maximises the average attention resource, and this is precisely the optimisation meaning of the `attention cycle'.
\end{example}

\section{The emotional cycle}

\begin{definition}[The circumplex model of emotion]\label{def:emotion_circumplex}
An emotional state can be represented on a two-dimensional plane: the horizontal axis is \textbf{emotional valence} (pleasant--unpleasant) and the vertical axis is \textbf{arousal level} (high arousal--low arousal). A particular emotion corresponds to a point in the plane: happiness (high pleasantness + high arousal), sadness (low pleasantness + low arousal), anxiety (low pleasantness + high arousal) and calm (high pleasantness + low arousal).

Over the course of a day the emotional state \textbf{rotates periodically} about the centre of the circumplex model---from high arousal in the morning (whether pleasant or not) to low arousal in the afternoon, then to low mood in the evening, and finally to a reset during sleep. The circadian variation of the emotional state can be regarded as a quasi-periodic trajectory in the (valence, arousal) plane: the emotional changes within a day correspond to an arc in the plane, and sleep resets the emotional state to near the origin (neutrality).
\end{definition}

\begin{proposition}[The emotional cycle of bipolar disorder]\label{prop:bipolar_cycle}
Bipolar disorder is characterised by periodic alternation of the emotional state between mania (high arousal + high pleasantness or high irritability) and depression (low arousal + low pleasantness). The time scale of one complete emotional cycle (mania--normal--depression--normal--mania) varies from person to person---about weeks to months for the rapid-cycling type and about several days for the ultra-rapid-cycling type.

In the phase space (valence $V$, arousal $A$), the emotional trajectory of a healthy individual is an ellipse of small amplitude (emotion fluctuates within the normal range), whilst the emotional trajectory of a patient with bipolar disorder is a limit cycle of large amplitude (emotion oscillates widely between the two extremes of mania and depression). The action of mood stabilisers (such as lithium salts) is to increase the rate of emotional return---constraining the orbit in phase space to a neighbourhood of the equilibrium point.
\end{proposition}

\begin{proof}
Linearise the emotional dynamics near the equilibrium point. Let $X = (V, A)^{\mathsf T}$ satisfy $\dot X = J X$, where the eigenvalues of $J$ are $\lambda_{\pm} = -\mu \pm i\omega$ ($\mu > 0$ is the rate of emotional return and $\omega$ the intrinsic oscillation frequency of emotion). For a healthy individual $\mu$ is relatively large: the orbit is a rapidly decaying spiral and the amplitude $r(t) = r_0 e^{-\mu t}$ is small. For a patient with bipolar disorder $\mu$ is markedly smaller (emotional return is weakened): the orbit is approximately closed and the limit-cycle amplitude $r \approx r_0$ is large. The action of mood stabilisers is to increase the effective $\mu$ (equivalent to introducing a stronger damping term in $J$), so that the orbit decays back into a neighbourhood of the equilibrium point.
\end{proof}

\section{The cycles of learning and forgetting}

\begin{definition}[The Ebbinghaus forgetting curve]\label{def:ebbinghaus}
The Ebbinghaus forgetting curve describes the exponential decay of memory retention over time in the absence of review:
\[
M(t) = M_0 \cdot e^{-t/\tau},
\]
where $M_0$ is the initial memory strength and $\tau$ the forgetting time constant (which depends on the difficulty of the memorised material and on the memory ability of the individual). Memory retention falls rapidly during the first few hours, and the rate of decline then slows.

The exponential decay of memory strength is analogous to the exponential decay of the Fourier components of the heat equation in the `Differential Equations and Analysis' chapter. The time constant $\tau$ measures the retention capacity of a memory trace under time translation. The larger $\tau$ is, the more lasting the memory.
\end{definition}

\begin{proposition}[The optimal review cycle for spaced repetition]\label{prop:spaced_repetition}
Spaced repetition exploits the periodic structure of the forgetting curve to optimise long-term memory. If each review resets the memory strength to $M_0$, then the optimal time interval before the $n$-th review is
\[
\Delta t_n = \tau \cdot \ln\frac{M_0}{M_{\min}} + \sum_{k=1}^{n-1} \Delta t_k \cdot \frac{\tau}{\tau + \Delta t_k} \approx \tau \cdot \ln\frac{M_0}{M_{\min}} \cdot n,
\]
where $M_{\min}$ is the minimum acceptable memory strength. The review interval grows linearly with the number of reviews---this is the scaling feature of the memory cycle.

Spaced repetition is the alternation, along the time axis, of the decay of memory and the reinforcement of review---each review resets the memory strength to its initial state ($M = M_0$), and the forgetting between two reviews constitutes one complete decay--recovery cycle. Optimising the review interval is equivalent to finding the cycle that maximises the average memory strength.
\end{proposition}

\begin{proof}
Suppose that after the $n$-th review the memory strength is $M_0$, and that after $\Delta t_n$ it has decayed to $M_{\min}$: $M_{\min} = M_0 e^{-\Delta t_n/\tau}$, which gives $\Delta t_n = \tau \ln(M_0/M_{\min})$. Consider the cumulative effect of review: if a review not only resets the memory strength but also slightly increases the forgetting time constant ($\tau_n = \tau_{n-1} + \Delta\tau$), then the review intervals grow gradually. Empirically, spaced-repetition software (such as Anki) uses a similar formula to compute optimal review intervals.
\end{proof}

\section{Subconscious processes}

\begin{definition}[Subconscious processes]\label{def:subconscious_process}
\textbf{Subconscious processes} are psychological processes that run automatically below the threshold of consciousness. Compared with conscious processes, subconscious processes have the following features:
\begin{enumerate}[label=(\roman*)]
    \item \textbf{Automaticity}: they can be initiated and run without conscious effort.
    \item \textbf{Non-introspectability}: their content cannot be observed directly by introspection.
    \item \textbf{Parallelism}: several subconscious processes can run simultaneously, unrestricted by the serial bottleneck of conscious attention.
    \item \textbf{Diversity of time scales}: they are distributed from the millisecond scale (implicit priming) through the hour scale (memory consolidation during sleep) to the interannual scale (the formation of implicit attitudes).
\end{enumerate}
The temporal structure of subconscious processes is an important component of psychological periodic structure---it provides the background rhythm of conscious processes and, through coupling with conscious processes, influences behavioural output.
\end{definition}

\begin{definition}[The activation--decay cycle of implicit memory]\label{def:implicit_memory}
\textbf{Implicit memory} (procedural memory, priming effects) manifests as memory traces that influence behaviour without conscious recall. The activation--decay of the implicit memory strength $I(t)$ over time can be modelled as
\[
\frac{dI}{dt} = -\lambda I + \eta \cdot \mathbf{1}_{\text{stimulus}}(t),
\]
where $\lambda$ is the decay rate and $\eta$ the activation gain. The decay rate $\lambda$ of implicit memory is usually smaller than the decay rate of explicit memory---this explains why procedural skills (such as riding a bicycle) are retained for longer than declarative knowledge (such as reciting poetry).
\end{definition}

\begin{proposition}[Coupling between subconscious and conscious cycles]\label{prop:subconscious_conscious_coupling}
Let the subconscious process $U(t)$ and the conscious process $C(t)$ interact through the following coupled equations:
\[
\frac{dU}{dt} = f_U(U, C), \qquad \frac{dC}{dt} = f_C(U, C).
\]
When the time scale of the subconscious process is close to that of the conscious process ($|\tau_U - \tau_C| < \epsilon$), \textbf{phase coupling} arises between the two processes---fluctuations in the subconscious state can influence the conscious state through the coupling, and vice versa.
\end{proposition}

\begin{proof}
Linearise the coupled equations near the equilibrium point: $\delta \dot U = a_{11}\delta U + a_{12}\delta C$, $\delta \dot C = a_{21}\delta U + a_{22}\delta C$. The eigenvalues of the coupled system are
\[
\lambda_{\pm} = \frac{(a_{11}+a_{22}) \pm \sqrt{(a_{11}-a_{22})^2 + 4 a_{12} a_{21}}}{2}.
\]
When the discriminant is negative, the eigenvalues are complex conjugates---the two processes exhibit coupled oscillation. The oscillation frequency is determined jointly by the coupling strength $a_{12}a_{21}$ and the intrinsic frequency difference $(a_{11}-a_{22})^2$. Under strong coupling the system enters a coupled periodic state---the subconscious and conscious processes oscillate in synchrony.
\end{proof}

\begin{example}[Memory consolidation during sleep]\label{ex:sleep_memory_consolidation}
Memory consolidation during sleep is a typical example of subconscious processes. Information learned during the day is first encoded in the hippocampus (a fast learning system), and during sleep the coupling between the hippocampus and the cortex (a slow learning system) gradually transfers the information to cortical long-term memory. This process takes place during the slow-wave sleep (N3) stage---at which the hippocampal--cortical coupling reaches its strongest.

From the perspective of periodic mathematics, the periodic structure of memory consolidation is nested within the sleep cycle: each N3 stage corresponds to one round of a `training cycle' of hippocampal--cortical coupling. Integrating the cumulative effect of several sleep cycles, the memory strength gradually increases over a whole night's sleep---this is an accumulative process driven by nested periods.
\end{example}

\section{Behavioural periodicity}

\begin{definition}[Behavioural period]\label{def:behaviour_periodicity}
A \textbf{behavioural period} is a repetitive behavioural pattern displayed by an individual on a particular time scale. Behavioural periods fall into the following levels:
\begin{enumerate}[label=(\roman*)]
    \item \textbf{Circadian behavioural rhythm} (about 24 hours): the circadian fluctuation of behavioural and physiological indices such as sleep--wake, eating, activity level, body temperature and hormone secretion.
    \item \textbf{Weekly behavioural period} (about 7 days): switching of activity patterns between workdays and rest days.
    \item \textbf{Habit loop} (minutes to hours): the three-step cue--behaviour--reward loop, driven by basal ganglia circuits.
    \item \textbf{Addiction cycle} (hours to days): the seeking--use--withdrawal cycle of a substance or behaviour.
    \item \textbf{Interannual behavioural period} (about 1 year): seasonal mood variation, festive behavioural patterns and seasonal changes in activity level.
\end{enumerate}
Behavioural periods at the various levels are nested within one another---the circadian behavioural rhythm contains several habit loops, the habit loop repeats within a week and is modulated by the weekly behavioural period, and the weekly behavioural period repeats within the interannual behavioural period.
\end{definition}

\begin{definition}[The three-stage model of the habit loop]\label{def:habit_loop}
The \textbf{habit loop} consists of three stages:
\begin{enumerate}[label=(\roman*)]
    \item \textbf{Cue}: an environmental stimulus that triggers the initial signal of the behaviour.
    \item \textbf{Routine}: the automatic sequence of behaviour triggered by the cue.
    \item \textbf{Reward}: the pleasure or satisfaction obtained after the behaviour is completed, which reinforces the cue--behaviour association.
\end{enumerate}
The period of the habit loop $T_{\text{habit}}$ is determined by the frequency of the cue and the strength of the reward. The higher the degree of automaticity of the habit, the shorter $T_{\text{habit}}$ is and the less the behaviour depends on conscious control.
\end{definition}

\begin{proposition}[Periodic characteristics of habit formation]\label{prop:habit_formation}
Let the behavioural tendency $B(t)$ vary with time; the process of habit formation can then be described by the following model:
\[
\frac{dB}{dt} = \gamma \cdot (B_{\infty} - B) \cdot \mathbf{1}_{\text{cue}}(t),
\]
where $\gamma$ is the learning rate and $B_{\infty}$ the asymptotic habit strength. The time scale of habit formation $\tau_{\text{habit}} = 1/\gamma$ is usually tens to hundreds of days---depending on the complexity of the behaviour and the frequency of the reward.
\end{proposition}

\begin{proof}
Under repeated cues, $B$ converges exponentially to $B_\infty$: $B(t) = B_\infty (1 - e^{-\gamma t})$. Experimental studies show that the average time of habit formation ranges from about 20 days for simple behaviours (such as drinking water) to about 250 days for complex behaviours (such as daily exercise). The time scale of habit formation $\tau_{\text{habit}}$ agrees with the rate of change of the neural markers of behavioural automaticity (the stability of basal ganglia activation).
\end{proof}

\begin{example}[The periodic structure of addiction]\label{ex:addiction_cycle}
Addictive behaviour exhibits a typical periodic structure:
\begin{enumerate}[label=(\roman*)]
    \item \textbf{Use phase}: the stimulus of the substance or behaviour activates the reward circuit (dopamine release), producing brief pleasure.
    \item \textbf{Tolerance phase}: repeated use reduces the sensitivity of the reward circuit (receptor downregulation), so that a stronger stimulus is required to produce the same effect.
    \item \textbf{Withdrawal phase}: after discontinuation the reward circuit is in a state of low activity---producing withdrawal symptoms such as anxiety, depression and craving.
    \item \textbf{Relapse phase}: the substance or behaviour is used again to relieve the withdrawal symptoms, forming a new round of the cycle.
\end{enumerate}
The time scale of the addiction cycle runs from hours (acute withdrawal) through days (prolonged withdrawal) to years (risk of relapse). The neural mechanisms of addiction involve long-term plasticity changes in the reward circuit---these changes `fix' the addiction cycle into a stable limit cycle, so that even if the external stimulus is absent for a long time, the internal craving mechanism may still drive relapse.
\end{example}

\section{Phase locking of psychological cycles}

\begin{definition}[Interpersonal emotional synchrony]\label{def:emotional_synchrony}
In face-to-face communication, the emotional states of the two interlocutors tend to \textbf{synchronise}---one person's pleasure or sadness is transmitted to the other through facial expression, tone of voice and body language. This emotional synchrony is the realisation in interpersonal interaction of the phase locking of the `Periodic Systems' chapter (Proposition~\ref{prop:phase-locking}).

Let the emotional valences of two people be $V_1(t)$ and $V_2(t)$, with intrinsic frequencies $\omega_1$ and $\omega_2$ respectively. Through emotional expression and perception the two become coupled:
\[
\frac{d\theta_1}{dt} = \omega_1 + c_{12} \sin(\theta_2 - \theta_1), \quad \frac{d\theta_2}{dt} = \omega_2 + c_{21} \sin(\theta_1 - \theta_2),
\]
where $\theta_i$ is the phase of the emotional state and $c_{ij}$ the emotional coupling strength (which depends on empathic ability and expressive intensity). When the coupling is strong enough, the emotional phases of the two people tend to coincide---emotional synchrony. Emotional synchrony between infant and caregiver is crucial to the development of a healthy attachment relationship.
\end{definition}

\begin{proposition}[The contagious spread of group emotion]\label{prop:emotional_contagion}
In a group setting, emotional states spread through the social network---one person's emotional change can influence others in their social neighbourhood, and thence individuals further away. This is analogous to periodic propagation on discrete lattice points in the `Periodic Manifold' chapter (Definition~\ref{def:lattice_periodic_evolution}).

Let the emotional valence $V_i(t)$ of individual $i$ in a group satisfy
\[
\frac{dV_i}{dt} = -\alpha V_i + \sum_{j \in \text{neighbourhood}(i)} w_{ij} (V_j - V_i) + \xi_i(t),
\]
where $\alpha$ is the rate of emotional return to neutrality, $w_{ij}$ the social influence weight and $\xi_i(t)$ an individual-specific emotional perturbation. When the coupling is strong ($w_{ij}$ large), the mean group emotion $\bar{V}(t) = \frac{1}{N}\sum_i V_i(t)$ exhibits a slow collective oscillation---the `social rhythm' of group emotion. The time scale of social emotional cycles (such as the economic optimism--pessimism cycle) can reach months to years---this is the macroscopic manifestation of psychological cycles at the social level.
\end{proposition}

\begin{proof}
Sum over all $i$ and divide by $N$: $\frac{d\bar V}{dt} = -\alpha \bar V + \frac{1}{N}\sum_{i} \sum_{j \in N(i)} w_{ij}(V_j - V_i) + \bar\xi$. When the weights are symmetric ($w_{ij} = w_{ji}$), the inner double sum cancels edge by edge in an undirected network: each edge $(i,j)$ contributes $w_{ij}(V_j - V_i) + w_{ji}(V_i - V_j) = 0$, so $\frac{d\bar V}{dt} = -\alpha \bar V + \bar\xi$. The mean therefore obeys a first-order linear equation---the collective oscillation does not come from the coupling term itself, but from the red-spectrum fluctuations of the noise $\bar\xi$ after low-pass filtering (time constant $1/\alpha$); the role of strong coupling is to align each $V_i$ rapidly with $\bar V$ (phase locking), thereby aggregating individual fluctuations into an observable social rhythm.
\end{proof}

\section{Conservation laws of psychological systems}

\begin{remark}[The three conservation laws of psychological systems (empirical observation / research programme)]\label{thm:psychological_conservation}
The dynamics of psychological systems may be described by three conservation laws. \textbf{The following three items are empirical observations and a research programme, not proved theorems}:
\begin{enumerate}[label=(\roman*)]
    \item \textbf{Energy conservation}: psychological energy (mental fatigue and recovery) is consumed during waking and restored during sleep. Sleep deprivation leads to a decline in cognitive function, and recovery sleep restores cognitive function---psychological energy satisfies conservation between consumption and recovery. This is the balance of energy under time translation: $\Delta E_{\text{consumed}} + \Delta E_{\text{recovered}} = 0$ (long-term average).
    \item \textbf{Information conservation}: the capacity of working memory is limited ($7 \pm 2$ chunks)---the inflow of information into working memory (encoding) and its outflow (forgetting or transfer to long-term memory) satisfy conservation. The consolidation of long-term memory takes place during sleep---the transfer of memory from the hippocampus to the cortex is a redistribution of information among different stores. The total information content of memory is conserved in the transfer.
    \item \textbf{Attention conservation}: the total attention resource is finite at any given moment---concentration of attention on one object is necessarily accompanied by neglect of other objects. The `zero-sum' character of attention manifests as the finiteness of the attention resource: the allocation of attention among several targets constitutes a phase constraint in phase space---focusing on one target means phase locking in the direction of that target.
\end{enumerate}
The hierarchical relation among the three conservation laws is: energy conservation guarantees the continuity of psychological activity, information conservation guarantees the stability of cognitive structure, and attention conservation guarantees the selectivity of the focus of consciousness. Together the three constitute the physical basis of the periodic structure of psychological systems.

\textbf{Empirical basis}:
(i) Energy conservation: the empirical evidence for the consumption and recovery of psychological energy comes from sleep-deprivation studies. Sustained waking leads to a progressive decline in cognitive function (reaction time, working memory, decision-making ability), and after recovery sleep cognitive function returns to the baseline level. The power of slow-wave activity (SWA, $0.5$--$4$ Hz) in the electroencephalogram (EEG) accumulates linearly during the waking period and decays exponentially during the sleep period---this is the electrophysiological index of psychological energy.

(ii) Information conservation: the limited capacity of working memory was established by the classical study of Miller (1956)---humans can hold about $7 \pm 2$ chunks of information simultaneously. Information is retained in working memory for about $20$--$30$ seconds, unless it is refreshed by rehearsal. The two-process theory of memory consolidation holds that the fast learning system of the hippocampus transfers information to the slow learning system of the cortex during sleep---the amount of information is conserved in the transfer.

(iii) Attention conservation: the limited-resource theory of attention holds that the total attention resource is fixed and finite at any given moment. Allocating attention to one task necessarily reduces the resources available for other tasks---dual-task interference experiments confirm the finiteness of the attention resource. In periodic mathematics, the attention resource $R(t)$ satisfies $R_1 + R_2 + \cdots + R_n \le R_{\max}$---total attention does not exceed the maximum capacity.
\end{remark}

\begin{remark}[The place of subconscious and behavioural periods in the book]
Subconscious processes and behavioural periodicity are two important complementary levels of the periodic structure of psychological systems. Subconscious processes provide a background rhythm beyond conscious processes---although these rhythms cannot be directly introspected, they can be observed indirectly through behavioural output, physiological indices and emotional fluctuation. Behavioural periodicity, in turn, projects the periodic structure of psychological systems onto the observable behavioural level---giving psychological cycles a measurable and verifiable empirical basis. Integrating the three levels of subconscious, conscious and behavioural processes within a single framework of periodic structure is the complete statement of this chapter's application of periodic mathematics in psychology.
\end{remark}

\begin{remark}[Connection with the rest of the book]
The theory of the periodic structure of psychological systems established in this chapter is an application of periodic mathematics in psychology. The sleep--wake cycle appears as a nested periodic structure formed by a main period of about 24 hours and an ultradian rhythm of about 90 minutes. The periodic fluctuation of attention appears as an alternation of focusing and dispersion. The circadian rotation in the circumplex model of emotion appears as epicyclic motion in the (valence, arousal) phase space. The Ebbinghaus forgetting curve appears as the exponential decay of memory strength, and spaced repetition optimises the cycle of memory retention. Subconscious processes provide a background rhythm beyond conscious processes---implicit memory, habituation and memory consolidation during sleep constitute the temporal structure of the subconscious. Behavioural periodicity projects psychological cycles onto the observable behavioural level---circadian behavioural rhythm, the weekly behavioural period, the habit loop and the addiction cycle constitute the multilevel periods of behaviour. Interpersonal emotional synchrony and group emotional contagion are the realisation, at the social-psychological level, of the phase locking of the `Periodic Systems' chapter. The three conservation laws of psychological systems---energy conservation, information conservation and attention conservation---together constitute the guarantee of self-consistency for the periodic structure of psychological systems.
\end{remark}

\section*{Seven-theme comparison table for this chapter}

\begin{center}
\begin{tabular}{>{\centering\arraybackslash}p{2.2cm}>{\centering\arraybackslash}p{10.5cm}}
\toprule
Theme & Corresponding content in this chapter \\
\midrule
Translation & The approximately 24-hour main period of sleep--wake: gated triggering by the two-process model (Proposition~\ref{prop:two_process_sleep}); linear accumulation of sleep pressure during the waking period \\
Scaling & Exponential decay of the Ebbinghaus forgetting curve $M(t) = M_0 e^{-t/\tau}$ (Definition~\ref{def:ebbinghaus}); scaling optimisation of spaced repetition (Proposition~\ref{prop:spaced_repetition}); multilevel nesting of behavioural periods (Definition~\ref{def:behaviour_periodicity}) \\
Rotation & Circadian rotation in the two-dimensional valence--arousal phase space (Definition~\ref{def:emotion_circumplex}); the large-amplitude limit cycle of bipolar disorder (Proposition~\ref{prop:bipolar_cycle}); coupled oscillation of subconscious and conscious processes (Proposition~\ref{prop:subconscious_conscious_coupling}) \\
Conservation & Psychological energy conservation, information conservation (working memory of $7\pm2$ chunks) and attention conservation ($R_1+\cdots+R_n \le R_{\max}$) (empirical observation / research programme, Remark~\ref{thm:psychological_conservation}) \\
Symmetry & Positive--negative valence symmetry of emotion ($V \in [-1,1]$, symmetric about $0$); circadian symmetry of sleep--wake \\
Closure & The consumption--recovery charge--discharge cycle of the attention resource (Lemma~\ref{ch26_attention_cycle}, Example~\ref{ch26_pomodoro_example}); the sleep reset of the emotional state; the cue--behaviour--reward closure of the habit loop (Definition~\ref{def:habit_loop}); the addiction cycle (Example~\ref{ex:addiction_cycle}) \\
Period & The 90-minute sleep cycle (NREM--REM alternation, Definition~\ref{def:sleep_cycle}); the theta ($4$--$8$ Hz) and delta ($0.5$--$2$ Hz) rhythms of attention oscillation (Definition~\ref{def:attention_oscillation}); phase locking of interpersonal emotional synchrony (Definition~\ref{def:emotional_synchrony}); the periodic characteristics of habit formation (Proposition~\ref{prop:habit_formation}); the activation--decay cycle of subconscious processes (Definition~\ref{def:implicit_memory}) \\
\bottomrule
\end{tabular}
\end{center}

\section*{Exercises}

\begin{enumerate}[leftmargin=3em, label=\textbf{Exercise \arabic*}]
    \item Prove the \textbf{periodicity decision theorem for psychological systems} of this chapter (\ref{thm:psych_periodicity_decision}). Hint: first write down the decision map and the evolution group of this chapter's target, then apply the proof structure of the general criterion given in the `Periodic Systems' chapter.

    \item Verify that the hypotheses of the \textbf{monotonicity of cardinality adjustment} (\ref{prop:psych_adaptive_growth}) hold for the domain objects of this chapter, and give a counterexample in which the hypotheses fail.

    \item Examine the behaviour of the \textbf{three-parameter correspondence: radius, phase and speed} (\ref{prop:psych_epicycle_correspondence}) under parameter degeneracy, and explain how the conclusion degenerates accordingly.

    \item Explain the relation between the \textbf{psychological periodicity diagnostic pipeline and its interface with three chapters} (\ref{thm:psych_diagnosis_integration}) and the corresponding general conclusion of the `Periodic Systems' chapter (whether it is identical, a special case, or requires additional conditions).

    \item Using the seven-theme comparison table of this chapter, identify the concrete carriers of the themes \textbf{translation}, \textbf{scaling} and \textbf{rotation} in this chapter, and give one example of domain data for each.

    \item Synthesis: restate the diagnostic pipeline of this chapter using the five-step standard interface of the `Periodic Systems' chapter (Table~\ref{tab:ps_five_steps}), noting one by one the theorem of this chapter on which each step depends: \ref{thm:psych_periodicity_decision}, \ref{prop:psych_adaptive_growth}, \ref{prop:psych_epicycle_correspondence}, \ref{thm:psych_diagnosis_integration}.
\end{enumerate}


